\documentclass[11pt]{article}
\usepackage{setspace}
\usepackage{geometry}
\geometry{margin=0.85in}
\setstretch{1.0}

\begin{document}

\begin{flushright}
Zhongli Wang \\
+49 160 9541 6015 \\
nkwangzhongli@gmail.com \\
Bielefeld, Deutschland \\
\end{flushright}

\vspace{0.1cm}

\noindent
EY-Parthenon \\
Recruiting Team \\
Eschborn / M\"unchen \\

\vspace{0.1cm}

\noindent
Bewerbung als Consultant Valuation, Modelling \& Economics -- Financial Services \\

\vspace{0.05cm}

\noindent Sehr geehrte Damen und Herren,

mit gro{\ss}em Interesse bewerbe ich mich f\"ur diese Position. Besonders reizvoll finde ich die Verbindung aus strategischer Beratung, transaktionsbezogener Unternehmensbewertung und datenbasierter Analyse in der Finanzindustrie. Die M\"oglichkeit, Banken, Versicherungen, und FinTech-Unternehmen bei Bewertungsfragen, strategischen Handlungsoptionen und Transformationen zu unterst\"utzen, entspricht sehr gut meinem Interesse an finanzwirtschaftlichen Fragestellungen und quantitativer Entscheidungsunterst\"utzung.

Derzeit schlie{\ss}e ich meine Promotion in Economics an der Universit\"at Bielefeld und der Universit\`a Cattolica del Sacro Cuore ab; der Abschluss ist f\"ur September 2026 vorgesehen. In meiner Forschung besch\"aftige ich mich mit komplexen Wirtschafts- und Unternehmenssystemen, datengetriebener Modellierung und statistischer Analyse. Dabei habe ich umfangreiche Erfahrung darin gesammelt, Probleme zu strukturieren, relevante Daten aufzubereiten, Annahmen kritisch zu hinterfragen und quantitative Ergebnisse klar zu kommunizieren. Diese F\"ahigkeiten m\"ochte ich nun in einem praxisnahen Beratungsumfeld einsetzen, in dem Analysen direkten Einfluss auf unternehmerische Entscheidungen haben.

Ein besonderer Schwerpunkt meiner bisherigen Arbeit liegt auf financial data analytics und Modellierung. In einem Forschungsprojekt zu Corporate-Fraud-Detection habe ich mit \"uberwachten Machine-Learning-Verfahren strukturierte Auff\"alligkeiten in Unternehmensdaten untersucht. Dadurch konnte ich meine Kenntnisse in Datenaufbereitung, Feature Engineering, Modellvalidierung und Interpretation quantitativer Ergebnisse vertiefen. In einem weiteren Dissertationsprojekt analysiere ich Finanz- und operative Daten von AT\&T und T-Mobile mithilfe simulationsbasierter Sch\"atzverfahren und neuronaler Surrogatmodelle. Diese Erfahrungen sind aus meiner Sicht gut anschlussf\"ahig an Bewertungs-, Modellierungs- und Analyseaufgaben im Financial-Services-Umfeld.
Dar\"uber hinaus bringe ich erste praktische Einblicke in das Bankgesch\"aft mit. 
W\"ahrend eines Praktikums bei der Agricultural Bank of China, Gansu Branch, konnte ich grundlegende Abl\"aufe einer kommerziellen Bank kennenlernen und ein erstes Verst\"andnis f\"ur bankinterne Prozesse, Abteilungsfunktionen und Koordination entwickeln. 

Neben meinen quantitativen F\"ahigkeiten bringe ich eine internationale und teamorientierte Arbeitsweise mit. Durch meine Promotion in einem europ\"aischen Forschungsumfeld, internationale Konferenzbeitr\"age und die Zusammenarbeit mit Betreuern und Forschenden aus verschiedenen L\"andern habe ich gelernt, komplexe Inhalte adressatengerecht aufzubereiten und konstruktiv im Team zu arbeiten. Ich verf\"uge \"uber sehr gute Deutsch- und Englischkenntnisse und kann mich schnell in neue Themen, Datenquellen und Projektkontexte einarbeiten.

EY-Parthenon \"uberzeugt mich, weil hier strategisches Denken, Transaktionserfahrung, Implementierungsn\"ahe und datenbasierte Analyse zusammenkommen. \"Uber die Gelegenheit, meine Motivation und meine Passung f\"ur diese Position in einem pers\"onlichen Gespr\"ach n\"aher zu erl\"autern, w\"urde ich mich sehr freuen.

\vspace{0.15cm}

\noindent Mit freundlichen Gr\"u{\ss}en \\[0.5cm]
\noindent Zhongli Wang

\end{document}
